\documentclass[11pt]{article}
\usepackage{amsmath,amssymb,amsthm}
\usepackage[margin=1in]{geometry}

\newtheorem{theorem}{Theorem}
\newtheorem{lemma}{Lemma}
\newtheorem{proposition}{Proposition}
\newcommand{\bits}{\{0,1\}}
\newcommand{\sat}{\#\mathrm{SAT}}

\title{Laboratory V94: Fixed-Language Child-Count Comparison}
\author{NC0\_k-Avoid Laboratory}
\date{2026-08-10}

\begin{document}
\maketitle

\paragraph{Prefix counts.}
For $C:\bits^N\to\bits^M$ and $p\in\bits^j$, write
\[
 N_C(pb)=\bigl|\{x:C_{<j}(x)=p,\ C_j(x)=b\}\bigr|.
\]
The canonical V92 bit is zero exactly when $N_C(p0)\le N_C(p1)$.

\paragraph{Fixed language.}
For $\alpha,\beta,c\in\bits$, let
\[
 D_{\alpha,\beta}(z,a,b)=1
 \Longleftrightarrow
 z=((a\oplus\alpha)\lor(b\oplus\beta)),
\]
and
\[
 Q_{c,\alpha}(s,u,z)=1
 \Longleftrightarrow
 (s\ne c)\lor(u\oplus\alpha)\lor z.
\]
Together with the projection $P(s)=s$, these are nine fixed gate types of
maximum arity three; call the language $\Gamma_{94}$.

\begin{theorem}[Arbitrary-prefix comparison]
The problem of deciding $N_C(p0)\le N_C(p1)$ for an arbitrary supplied prefix
$p$ is PP-complete for $\Gamma_{94}$ circuits even under the exact stretch
restriction $M=N+1$.
\end{theorem}

\begin{proof}
Membership follows because the two child counts are \#P functions and
$N_C(p1)-N_C(p0)+1$ is a GapP function whose strict positivity is exactly the
non-strict comparison.

For hardness, start from a pair of 3-CNF formulas $F_0,F_1$ over a common set
of $n$ variables.  Comparison of their numbers of satisfying assignments is
PP-complete: the standard GapP sign characterization gives \#P functions
$a,b$ with membership equivalent to $a>b$; compare instead the \#P functions
$2b+1$ and $2a$, and use a parsimonious, unique-extension Cook--Levin/Tseitin
encoding into width-three CNF.

Let $r$ be the total number of clauses in $F_0,F_1$.  Introduce a selector $s$,
one auxiliary $z_A$ for every source clause $A=(\ell_1\lor\ell_2\lor\ell_3)$,
and
\[
 D=\max\{0,r-n-1\}
\]
free dummy variables.  Thus
\[
 N=n+1+r+D.
\]
For every clause of $F_c$, place two prefix outputs: first a suitable signed
$D_{\alpha,\beta}$ gate enforcing
$z_A=\ell_2\lor\ell_3$, then a suitable $Q_{c,\alpha}$ gate enforcing
$(s\ne c)\lor\ell_1\lor z_A$.  Fix all $2r$ such outputs to one.  The next
output is $P(s)$, and harmless projections pad the circuit to $M=N+1$; this is
possible because $2r\le N$ by the definition of $D$.

For every assignment of the source variables, the definition gates determine
all auxiliaries uniquely.  On branch $s=c$, conditional gates from the other
formula are tautologies and those from $F_c$ enforce exactly its source
clauses.  The $D$ dummies remain free.  Consequently
\[
 N_C(p0)=2^D\sat(F_0),\qquad
 N_C(p1)=2^D\sat(F_1).
\]
The common factor preserves the comparison, proving PP-hardness.
\end{proof}

\begin{proposition}[Canonical-prefix separation]
For the explicit ordering above, the hard prefix is not the V92 canonical
prefix already at its first definition gate.
\end{proposition}

\begin{proof}
Every definition truth table has four accepting and four rejecting assignments
on its three local arguments.  With all other inputs free, the two first-child
counts are equal, so the V92 tie rule chooses zero.  The reduction fixes that
prefix output to one.
\end{proof}

\begin{theorem}[Affine control]
If every output of $C$ is affine over $\mathbb F_2$, exact child comparison is
polynomial-time for every prefix, and the V92 canonical avoided output for
$M=N+1$ is computable in polynomial time.
\end{theorem}

\begin{proof}
Every fixed output bit is one linear equation.  Gaussian elimination gives
consistency and rank, hence each child count is either zero or
$2^{N-\operatorname{rank}}$.  Incrementally, if the next coefficient row is
independent then the children have equal size and the canonical bit is zero.
If it is dependent, exactly one output bit is forced on the current nonempty
fiber, so the canonical rule selects the opposite empty child.  Once empty,
all remaining comparisons tie at zero.  With $N+1$ outputs an empty prefix is
therefore obtained, yielding a word outside the range.
\end{proof}

\paragraph{Nonclaims.}
The PP-completeness theorem concerns arbitrary supplied prefixes only.  It does
not prove PP-hardness of V92-generated canonical prefixes, does not prove that
range avoidance requires exact comparison, does not improve the
$O(n2^{n/2})$ all-instance $\mathrm{NC}^0_3$ range-avoidance baseline, and does
not resolve P versus NP.

\end{document}
